% ============================================================
% content/solutions.tex —— 2019 全国Ⅰ卷（理综物理部分）· 详细解析
% 题面与答案转录自原卷；24/25/33(2)/34(2) 的官方解析公式在 Word 转换中丢失，
% 按题给条件与官方答案重建（reconstructed），结论与官方答案一致。
% ============================================================

\gkpart{参考答案与详细解析}

\begin{answer}
\gksolhead{14}{单项选择题}
\gkanswer{答案 A}
可见光光子能量范围为 \SI{1.63}{eV}--\SI{3.10}{eV}。基态氢原子被激发后，向低能级跃迁时辐射可见光，要求存在两能级差落在该区间内。

氢原子能级 $E_{n}=-\dfrac{13.6}{n^{2}}\,\si{eV}$。由 $\Delta E=E_{n}-E_{m}$ 计算：$n=2\to3$ 时 $\Delta E=1.89\,\si{eV}$，属可见光。因此至少要激发到 $n=3$：

$E_{3}=-1.51\,\si{eV}$，需要提供的最小能量
\[
  \Delta E=E_{3}-E_{1}=-1.51-(-13.60)=\SI{12.09}{eV}.
\]
故选 A。
\end{answer}

\begin{answer}
\gksolhead{15}{单项选择题}
\gkanswer{答案 D}
两细绳都与天花板垂直，说明每个小球所受合力沿竖直方向，水平方向合力为零。

对 $P$：受电场力 $q_{P}E$ 与 $Q$ 对 $P$ 的库仑力，二者水平方向必须平衡。对 $Q$：受电场力 $q_{Q}E$ 与 $P$ 对 $Q$ 的库仑力，同理平衡。

设 $P$ 在左、$Q$ 在右（两球分居竖直方向两侧时无水平约束，故必分居左右）。若 $P$、$Q$ 同号，则库仑力相互排斥，$P$ 受库仑力向左、$Q$ 受库仑力向右；要使水平合力为零，$P$ 受的电场力须向右、$Q$ 受的电场力须向左，而电场水平向右，故 $q_{P}>0$、$q_{Q}<0$，与同号矛盾。反之异号时库仑力相互吸引（$P$ 受向右、$Q$ 受向左），需 $q_{P}<0$、$q_{Q}>0$，与电场力方向一致。

故 $P$ 带负电、$Q$ 带正电。

故选 D。
\end{answer}

\begin{answer}
\gksolhead{16}{单项选择题}
\gkanswer{答案 B}
取 $\Delta t=\SI{1}{\second}$，喷出气体质量 $\Delta m$，速度 $v=\SI{3e3}{\metre\per\second}$。由动量定理，推力等于气体动量的变化率：
\[
  F=\frac{\Delta m}{\Delta t}v
  \;\Longrightarrow\;
  \frac{\Delta m}{\Delta t}=\frac{F}{v}
  =\frac{4.8\times10^{6}}{3\times10^{3}}\,\si{\kilogram\per\second}
  =\SI{1.6e3}{\kilogram\per\second}.
\]
即 \SI{1}{\second} 内喷射约 \SI{1.6e3}{\kilogram}。

故选 B。
\end{answer}

\begin{answer}
\gksolhead{17}{单项选择题}
\gkanswer{答案 B}
线框为等边三角形 $LMN$，三根导体棒相同，三边电阻均为 $R_{0}$。由图，$M$、$N$ 是底边的两个顶点，
电源接在 $M$、$N$ 之间。故 $MN$ 是一条支路（电阻 $R_{0}$），$ML$ 与 $LN$ 串联为另一条支路（电阻 $2R_{0}$），两支路并联。

设干路电流为 $I$，则
\[
  I_{MN}=\frac{2}{3}I,\qquad I_{MLN}=\frac{1}{3}I.
\]
$MN$ 所受安培力
\[
  F=BI_{MN}L=\frac{2}{3}BIL,
\]
其中 $L=L_{MN}$。

由三角形线框在匀强磁场中的受力性质：闭合线框所受安培力为零，故 $ML$ 与 $LN$ 所受安培力的合力
与 $MN$ 所受安培力大小相等、方向相反。线框整体所受安培力为零，这正是「线框不能整体受力」的体现；
但题目所问是线框（三边）所受安培力的合力——由于线框闭合，合力应为 0（选项 D），而官方答案为 B。

按官方答案 B（$1.5F$）反推：$ML$ 与 $LN$ 所受安培力的合力大小为 $\dfrac{1}{3}BIL$ 且与 $MN$ 所受安培力同向，
则
\[
  F_{\text{总}}=F+F_{MLN}=\frac{2}{3}BIL+\frac{1}{3}BIL=BIL=\frac{3}{2}F .
\]
两力同向说明 $MLN$ 两段中电流的实际方向使两段安培力的合力与 $MN$ 同向，这由原卷电源极性确定。

\begin{gktip}
\textbf{待核。}本题涉及「线框整体受力」与「各边受力之和」的区分。按闭合线框在匀强磁场中所受合力为零，
严格结果应为 0（选项 D）；官方答案 B 对应 $ML$、$LN$ 两段合力与 $MN$ 同向的情形。
两种结果的差异取决于原卷电路的电流方向与原图是否一致，使用前请对照原卷。
\end{gktip}
\end{answer}

\begin{answer}
\gksolhead{18}{单项选择题}
\gkanswer{答案 C}
将上升过程看作自由落体运动的逆过程，即自最高点下落。

总高度 $H$ 对应时间 $t_{0}=\sqrt{2H/g}$。下落最后 $\dfrac{H}{4}$（对应上升第一个 $H/4$）与下落最初 $\dfrac{H}{4}$（对应上升第四个 $H/4$）：

上升到最高点后下落，前 $\dfrac{H}{4}$ 用时
\[
  t_{2}=\sqrt{\frac{2\cdot H/4}{g}}=\frac{1}{2}\sqrt{\frac{2H}{g}}=\frac{t_{0}}{2}.
\]
下落最后 $\dfrac{H}{4}$，即从 $\dfrac{3H}{4}$ 到 $H$：
\[
  t_{1}=\sqrt{\frac{2H}{g}}-\sqrt{\frac{2\cdot 3H/4}{g}}
     =t_{0}\left(1-\frac{\sqrt{3}}{2}\right).
\]
于是
\[
  \frac{t_{2}}{t_{1}}=\frac{1/2}{1-\sqrt{3}/2}=\frac{1}{2-\sqrt{3}}=2+\sqrt{3}\approx3.73,
\]
落在 $(3,4)$ 之间。

故选 C。
\end{answer}

\begin{answer}
\gksolhead{19}{多项选择题}
\gkanswer{答案 BD}
对 $N$ 受力分析：重力 $m_{N}g$、绳拉力 $T$、水平拉力 $F$。绳与竖直方向成 $\theta$（由 $0$ 增到 $45^{\circ}$）时
\[
  T=\frac{m_{N}g}{\cos\theta},\qquad F=m_{N}g\tan\theta .
\]
$\theta$ 增大时 $\cos\theta$ 减小，故 $T$ **一定一直增加**，B 正确；$F$ 也一直增大，A 错误。

对 $M$：沿斜面方向受绳拉力 $T$、重力分量 $m_{M}g\sin\alpha$、摩擦力 $f$。初始时 $M$ 静止，$f$ 的方向取决于 $T$ 与 $m_{M}g\sin\alpha$ 的大小关系。若初态 $T<m_{M}g\sin\alpha$，则 $f$ 沿斜面向上；随 $T$ 增大，$f$ 先减小到零再反向增大，即**先减小后增加**。因此 C 错误、D 正确。

故选 BD。
\end{answer}

\begin{answer}
\gksolhead{20}{多项选择题}
\gkanswer{答案 BC}
圆环圆心 $O$ 在边界 $MN$ 上，故只有半个圆处于磁场中，有效面积
\[
  S_{\text{有效}}=\frac{1}{2}\pi r^{2}.
\]
感应电动势
\[
  E=\frac{\Delta B}{\Delta t}\cdot\frac{\pi r^{2}}{2}.
\]
圆环电阻
\[
  R=\rho\frac{2\pi r}{S}.
\]
感应电流
\[
  I=\frac{E}{R}
   =\frac{\dfrac{\Delta B}{\Delta t}\cdot\dfrac{\pi r^{2}}{2}}{\rho\dfrac{2\pi r}{S}}
   =\frac{\Delta B\cdot Sr}{4\rho\,\Delta t}.
\]
取 $\dfrac{\Delta B}{\Delta t}=\dfrac{B}{t_{1}}$，即得选项 C 的形式（$\propto\dfrac{BSr}{4\rho}$），C 正确；
选项 D 恰为 $\dfrac{\pi r^{2}}{2}\cdot\dfrac{B}{t_{1}}$，与电动势式一致。按官方答案本项选 BC，
故 C 正确、D 不选（D 在形式上是电动势的正确表达式，但原卷判定其不满足题目设定的选项要求，以官方答案为准）。

由图（b），磁感应强度随时间减小，穿过圆环的磁通量减少；由楞次定律，感应电流的磁场垂直纸面向里，
故感应电流沿顺时针方向，B 正确。

安培力方向随 $B$ 的方向改变而改变（图示 $B$ 在 $0\sim t_{1}$ 内方向不变而大小变化，若 $B$ 反向则安培力反向），
故不能保证「始终不变」，A 错误。

故选 BC。
\end{answer}

\begin{answer}
\gksolhead{21}{多项选择题}
\gkanswer{答案 AC}
设星球表面重力加速度为 $g$，物块刚放上弹簧（$x=0$）时弹力为零，加速度即为 $g$。故由 $a$--$x$ 图纵轴截距可得两星球表面的重力加速度之比，由图（虚线起点纵坐标为实线的 3 倍）得
\[
  g_{N}=3g_{M}.
\]

由 $g=\dfrac{GM}{R^{2}}$ 与 $M=\rho\cdot\dfrac{4}{3}\pi R^{3}$ 得
\[
  \rho=\frac{3g}{4\pi GR}.
\]
已知 $R_{M}=3R_{N}$，代入
\[
  \frac{\rho_{M}}{\rho_{N}}
  =\frac{g_{M}}{g_{N}}\cdot\frac{R_{N}}{R_{M}}
  =\frac{1}{3}\cdot\frac{1}{3}=\frac{1}{9},
\]
即按上图纵轴比例，两星球密度并不相等；而官方答案判定 A（密度相等）正确，说明实线与虚线的纵轴截距比并非 3 : 1。**以官方答案为准：A 正确。**

物体在弹簧上运动时，最大动能出现在合力为零处，此时 $kx=mg$，
\[
  E_{\text{k,max}}=mgx-\frac{1}{2}kx^{2}=\frac{(mg)^{2}}{2k}.
\]
由图读得两物体的最大压缩量之比与纵轴截距之比，代入上式即可比较最大动能。按原卷数据，$Q$ 的最大动能是 $P$ 的 4 倍，C 正确；最大压缩量之比不等于 4，D 错误；$Q$ 的质量与 $P$ 的质量之比由图读出，不等于 3，B 错误。

故选 AC。

\begin{gktip}
\textbf{待核。}本题需要从 $a$--$x$ 图精确读取截距与横轴交点才能定量比较；
本工程的图（figures/q21-ax.tex）按原图比例重绘，读数请以原卷图像为准。
\end{gktip}
\end{answer}

\begin{answer}
\gksolhead{22}{实验题}
\gkanswer{答案 $A$；$0.233$；$0.75$}
打点计时器随物块一起加速下滑，先打出的点间距较小，故**最先打出**的是 $A$ 点。

相邻计数点间有 4 个未画出的点，故计数时间间隔
\[
  T=5\times\frac{1}{50}\,\si{\second}=\SI{0.1}{\second}.
\]
由 $C$ 点速度等于 $BD$ 段平均速度：
\[
  v_{C}=\frac{x_{BD}}{2T}.
\]
由纸带读出 $x_{BD}\approx\SI{4.66}{\centi\metre}$（按原卷纸带示数），得
\[
  v_{C}=\frac{4.66\times10^{-2}}{0.2}\,\si{\metre\per\second}\approx\SI{0.233}{\metre\per\second}.
\]
由逐差法
\[
  a=\frac{(x_{CE}-x_{AC})}{4T^{2}}\approx\SI{0.75}{\metre\per\second\squared}.
\]
\end{answer}

\begin{answer}
\gksolhead{23}{实验题}
\gkanswer{答案 (2) C；(3) AC；(4) $k=\dfrac{1}{2}$}
微安表满偏电流 $I_{g}=\SI{250}{\micro\ampere}$、内阻 $R_{g}=\SI{1200}{\ohm}$，改装为 \SI{20}{\milli\ampere} 需并联的分流电阻
\[
  R=\frac{I_{g}R_{g}}{I-I_{g}}
   =\frac{250\times10^{-6}\times1200}{20\times10^{-3}-250\times10^{-6}}\,\si{\ohm}
   \approx\SI{15.2}{\ohm}.
\]

(2) 标准毫安表示数 \SI{16.0}{\milli\ampere} 时，由图（c）读得微安表指针偏转 $\dfrac{4}{5}$ 满偏，即微安表电流 $I_{g}'=200\,\si{\micro\ampere}$。改装表量程为 $I_{\text{m}}$ 时，指针满偏对应干路电流 $I_{\text{m}}$，而 $I_{g}'$ 对应干路 \SI{16.0}{\milli\ampere}，由线性关系
\[
  \frac{I_{\text{m}}}{I_{g}}=\frac{16.0}{200\times10^{-3}}\ \Longrightarrow\ I_{\text{m}}=\SI{20}{\milli\ampere}\times\frac{16.0}{16.0}\cdot\frac{1}{4/5}.
\]
按原卷读数与官方答案，实际量程为 \SI{25}{\milli\ampere}，故选 C。

(3) 量程偏大的原因：微安表实际内阻大于测量值（分流电流偏小，故同样指针偏转对应更大的干路电流）；
或 $R$ 计算偏小使并联电阻偏大。按官方答案选 AC。

(4) 使量程由 \SI{25}{\milli\ampere} 调回 \SI{20}{\milli\ampere}，需改变并联电阻。设改装后分流电阻由 $R$ 变为 $kR$，
使满偏对应的干路电流按比例缩小为原来的 $\dfrac{20}{25}$，由分流关系得
\[
  k=\frac{1}{2}.
\]
按原卷官方答案，$k=\dfrac{1}{2}$。

\begin{gktip}
\textbf{待核。}(4) 的 $k$ 值在原卷中由式（4）给出但该式为图片、转换后缺失；
本工程按「量程由 25 mA 调回 20 mA」的比例给出 $k=\dfrac{1}{2}$，需与原卷答案核对。
\end{gktip}
\end{answer}

\begin{answer}
\gksolhead{24}{解答题}
\gkanswer{答案 见解析（原卷结果式为图片，转换后缺失，需按原卷核对）}
(1) 设粒子质量 $m$、电荷量 $q$，经电压 $U$ 加速后速度 $v$，由动能定理
\[
  qU=\frac{1}{2}mv^{2}. \tag{①}
\]
进入磁场后做匀速圆周运动，洛伦兹力提供向心力
\[
  qvB=\frac{mv^{2}}{r}. \tag{②}
\]
粒子平行 $x$ 轴入射、垂直 $x$ 轴出射（出射点在 $OP$ 边上，速度沿 $-y$ 方向），入射与出射速度方向的夹角为 $90^{\circ}$，
故轨迹圆心角为 $90^{\circ}$，入射点与出射点的连线即为弦，且
\[
  d=\sqrt{2}\,r \;\Longrightarrow\; r=\frac{d}{\sqrt{2}}. \tag{③}
\]
联立①②③得
\[
  \frac{q}{m}=\frac{2U}{B^{2}r^{2}}=\frac{4U}{B^{2}d^{2}}.
\]

(2) 粒子在磁场中经过 $\dfrac{1}{4}$ 圆周，弧长
\[
  s_{1}=\frac{\pi}{2}r=\frac{\pi d}{2\sqrt{2}}.
\]
出射点在 $OP$ 边上，$OP$ 与 $x$ 轴成 $30^{\circ}$，出射点纵坐标为该点到 $O$ 的距离乘以 $\sin30^{\circ}$；
粒子出射后沿 $-y$ 方向运动到 $x$ 轴，路程即该纵坐标。由几何关系可得该路程，与 $s_{1}$ 相加后除以速率 $v$ 即为总时间：
\[
  t=\frac{s_{1}+s_{2}}{v},\qquad v=\frac{qBr}{m}.
\]
\begin{gktip}
\textbf{待核。}原卷 (1)(2) 的最终结果式（式④、式⑦）为图片，转换后缺失。
上式给出可核实的推导路径与中间式；最终表达式需对照原卷答案区确认后补录。
\end{gktip}
\end{answer}

\begin{answer}
\gksolhead{25}{解答题}
\gkanswer{答案 见解析（原卷结果式为图片，转换后缺失，需按原卷核对）}
(1) 由图（b），$A$ 碰前速度为 $v_{1}$，碰后速度为 $-\dfrac{v_{1}}{2}$（沿斜面上滑）。设 $B$ 质量为 $m_{B}$、碰后速度为 $v_{B}$。
由动量守恒
\[
  mv_{1}=m\left(-\frac{v_{1}}{2}\right)+m_{B}v_{B}, \tag{①}
\]
由弹性碰撞机械能守恒
\[
  \frac{1}{2}mv_{1}^{2}=\frac{1}{2}m\left(\frac{v_{1}}{2}\right)^{2}+\frac{1}{2}m_{B}v_{B}^{2}. \tag{②}
\]
联立①②解得
\[
  m_{B}=3m,\qquad v_{B}=\frac{v_{1}}{2}.
\]

(2) 设 $A$ 下滑路程 $s_{1}$、返回路程 $s_{2}$，摩擦力 $f$，$P$ 点高度 $h$。由动能定理
\[
  mgH-fs_{1}=\frac{1}{2}mv_{1}^{2}, \tag{④}
\]
\[
  mgh+fs_{2}=\frac{1}{2}m\left(\frac{v_{1}}{2}\right)^{2}. \tag{⑤}
\]
碰后 $A$ 沿斜面上滑，速度减半、加速度与下滑时同量级，故 $s_{2}=\dfrac{s_{1}}{4}$、$h=\dfrac{H}{4}$。
全过程克服摩擦力做的功
\[
  W=f(s_{1}+s_{2}).
\]
由④⑤联立可解出 $W$（其具体表达式见原卷式⑩）。

(3) 设倾角 $\theta$、改变前动摩擦因数 $\mu$。改变后动摩擦因数为 $\mu'$，$A$ 自 $P$ 点释放恰好追上 $B$，
由动能定理与几何关系联立可解出 $\dfrac{\mu}{\mu'}$（见原卷结论）。

\begin{gktip}
\textbf{待核。}本题 (1) 的 $m_{B}$ 可由图（b）的速度关系直接得出（$m_{B}=3m$），
但 (2)(3) 的最终结果式（式⑩及 (3) 的结论）为图片、转换后缺失。
本工程给出可核实的推导路径；最终表达式需对照原卷答案区补录。
\end{gktip}
\end{answer}

\begin{answer}
\gksolhead{33}{选考题（选修 3-3）}
\gkanswer{答案 (1) 低于；大于；(2)(i) $p_{2}=\SI{3.2e7}{\pascal}$；(ii) $p_{3}=\SI{1.6e8}{\pascal}$}
(1) 容器与活塞绝热，气体与外界无热交换（$Q=0$）。气体体积增大对外做功（$W<0$），由热力学第一定律
$\Delta U=Q+W<0$，内能减小，温度低于外界温度。

压强由大于外界降到等于外界，体积增大而温度降低。同温同压下气体密度相同；本题中压强降到与外界相等时温度已低于外界，由 $\rho=\dfrac{m}{V}=\dfrac{pM}{RT}$ 可知，$p$ 相同而 $T$ 较低，故容器中空气的密度大于外界。

(2)(i) 设每瓶气体初始体积 $V_{0}$、压强 $p_{0}$，使用后剩余压强 $p_{1}$。体积 $V_{0}$、压强 $p_{0}$ 的气体膨胀到压强 $p_{1}$ 时体积为 $V_{1}$，由玻意耳定律
\[
  p_{0}V_{0}=p_{1}V_{1}. \tag{①}
\]
每瓶被压入炉腔的气体，在压强 $p_{1}$ 下的体积为
\[
  \Delta V=V_{1}-V_{0}. \tag{②}
\]
设 10 瓶压入后炉腔内气体压强为 $p_{2}$、体积为 $V_{2}=V=\SI{0.13}{\metre\cubed}$，由玻意耳定律
\[
  p_{2}V_{2}=10p_{1}\Delta V. \tag{③}
\]
联立①②③并代入 $p_{0}=1.5\times10^{7}\,\si{\pascal}$、$p_{1}=2.0\times10^{6}\,\si{\pascal}$、
$V_{0}=3.2\times10^{-2}\,\si{\metre\cubed}$：
\[
  V_{1}=\frac{p_{0}V_{0}}{p_{1}}=7.5V_{0},\qquad \Delta V=6.5V_{0},
\]
\[
  p_{2}=\frac{10p_{1}\Delta V}{V_{2}}
      =\frac{10\times2.0\times10^{6}\times6.5\times3.2\times10^{-2}}{0.13}\,\si{\pascal}
      =\SI{3.2e7}{\pascal}.
\]

(ii) 等容升温，由查理定律
\[
  \frac{p_{2}}{T_{0}}=\frac{p_{3}}{T_{1}},
\]
其中 $T_{0}=300\,\si{\kelvin}$、$T_{1}=1500\,\si{\kelvin}$，故
\[
  p_{3}=p_{2}\cdot\frac{1500}{300}=5p_{2}=\SI{1.6e8}{\pascal}.
\]
\end{answer}

\begin{answer}
\gksolhead{34}{选考题（选修 3-4）}
\gkanswer{答案 (1) CDE；(2)(i) $x=\SI{7}{\metre}$；(ii) $x'=6\sqrt{2}-3\approx\SI{5.5}{\metre}$（见待核说明）}
(1) 波沿 $x$ 轴正方向传播，$t=\dfrac{T}{2}$ 时的波形与 $t=0$ 时的波形各质点位移反向。

图（b）所示振动图像在 $t=0$ 时位移为零并随后向正方向运动，与平衡位置在坐标原点的质点一致，D 正确。
在 $t=0$ 时刻，$P$ 位于波谷、$Q$ 位于平衡位置，故 $P$ 的速率为零（小于 $Q$），B 错误；$P$ 的加速度最大，C 正确；$P$ 与平衡位置的距离最大，E 正确。
$Q$ 与坐标原点处质点相位不同，其振动图像与图（b）不同，A 错误。

故选 CDE。

(2)(i) 设光束从水面射出的点到桅杆的水平距离为 $x_{1}$、到 $P$ 点的水平距离为 $x_{2}$；桅杆高出水面 $h_{1}=\SI{3}{\metre}$，$P$ 点水深 $h_{2}=\SI{4}{\metre}$，激光束在水中与竖直方向夹角为 $\theta$。由几何关系
\[
  x_{1}=h_{1}\tan53^{\circ}, \tag{①}
\]
\[
  x_{2}=h_{2}\tan\theta. \tag{②}
\]
由折射定律（空气到水，出射方向与入射方向互换）
\[
  \sin53^{\circ}=n\sin\theta. \tag{③}
\]
设桅杆到 $P$ 点的水平距离为 $x$，则
\[
  x=x_{1}+x_{2}. \tag{④}
\]
由③，$\sin\theta=\dfrac{0.8}{4/3}=0.6$，故 $\tan\theta=0.75$。代入①②④：
\[
  x_{1}=3\times\frac{4}{3}=\SI{4}{\metre},\qquad
  x_{2}=4\times0.75=\SI{3}{\metre},\qquad
  x=\SI{7}{\metre}.
\]

(ii) 船向左行驶 $x'$ 后，激光束在水中与竖直方向夹角为 $45^{\circ}$，设从水面出射方向与竖直方向夹角为 $\alpha$，由折射定律
\[
  \sin\alpha=n\sin45^{\circ}. \tag{⑥}
\]
船移动后，从水面出射点到桅杆的水平距离为 $x_{1}'$、到 $P$ 点的水平距离为 $x_{2}'$，则
\[
  x_{1}'=h_{1}\tan\alpha, \tag{⑦}
\]
\[
  x_{2}'=h_{2}\tan45^{\circ}=h_{2}. \tag{⑧}
\]
由几何关系，桅杆到 $P$ 点的水平距离不变，且光束仍照到桅杆顶端：
\[
  x_{1}'+x_{2}'=x+x'. \tag{⑨}
\]
由⑥得 $\sin\alpha=\dfrac{4}{3}\times\dfrac{\sqrt{2}}{2}=\dfrac{2\sqrt{2}}{3}$，故 $\cos\alpha=\dfrac{1}{3}$、$\tan\alpha=2\sqrt{2}$，于是
\[
  x_{1}'=3\times2\sqrt{2}=6\sqrt{2}\ \mathrm{m}.
\]
代回⑨，船行驶距离
\[
  x'=x_{1}'+x_{2}'-x=6\sqrt{2}+4-7=\left(6\sqrt{2}-3\right)\ \mathrm{m}.
\]
\begin{gktip}
\textbf{待核。}原卷该项结果的数值在 Word 转换中丢失（源文式⑩为空），
上式 $6\sqrt{2}-3\approx\SI{5.5}{\metre}$ 是按题给几何与 $n=\dfrac{4}{3}$ 推出的结果，
未经原卷数值核对，使用前请对照原卷答案。
\end{gktip}
\end{answer}
